\documentclass[acmsmall,11pt,nonacm]{acmart}
\geometry{reset,letterpaper,margin=1in,headheight=14pt,footskip=24pt}
\setcopyright{none}
\settopmatter{printacmref=false,printccs=false,printfolios=true}
\setcitestyle{numbers,square,nosort,nocompress}
\renewcommand{\bibfont}{\normalsize}
\usepackage{mathtools}
\newcommand{\ind}{\mathbf{1}}
\newcommand{\disc}{\operatorname{disc}}
\newcommand{\Zk}{\mathbb{Z}_k}
\newcommand{\calB}{\mathcal{B}}
\newcommand{\calE}{\mathcal{E}}
\newcommand{\calC}{\mathcal{C}}
\newcommand{\Phiic}{\Phi_{i,c}}
\title{A Size-Independent Discrepancy Bound for Sparse Coverage Functions}
\author{Zijian Zeng}
\affiliation{\institution{UCSI University}\city{}\country{Malaysia}}
\email{1002266693@ucsiuniversity.edu.my}
\authorsaddresses{\normalsize Zijian Zeng, UCSI University,
  \href{mailto:1002266693@ucsiuniversity.edu.my}{1002266693@ucsiuniversity.edu.my}.}
\renewcommand{\shortauthors}{Zijian Zeng}

\begin{document}
\begin{abstract}
\normalsize
A coverage function counts the features reached by a set of items. Consider
any family of such functions in which each item belongs to at most $t$
feature occurrences across the entire family. We prove that, for every
integer $k\ge 2$, the items can be partitioned into $k$ color classes so
that every function differs by at most $4kt(kt-1)^2$ between any two
classes. The bound is independent of the number of items, functions,
features, and the sizes of the features. It answers affirmatively the
qualitative question of whether sparse coverage discrepancy admits a
polynomial bound in $t,k$ alone, motivated by the size-dependent bound of
Avila, Rohwedder and Wennmann (ESA 2026). Its dependence on $t,k$ is coarser
than theirs.

The proof first uses Hall's theorem to assign disjoint $k$-item witnesses
to all features of size at least $kt$. Restricting a witness to cyclic
color shifts makes every such feature exactly balanced under every allowed
choice. A graph on the resulting blocks has at most $(kt-1)^2$ independent
classes of small-feature interactions. Within one class, coverage
expectations are affine in the block-state probabilities, with maximum
column sum at most $kt$. Applying the known multiple-choice Beck--Fiala
rounding theorem class by class proves the bound. The construction is
deterministic and polynomial time for explicitly listed incidences,
including binary-encoded $k$ and $t$.
\end{abstract}
\maketitle
\thispagestyle{plain}
\pagestyle{plain}
\clearpage

\section{The result and its context}
\label{sec:intro}

Coverage valuations measure diversity: reaching the same feature twice
does not increase its contribution. This saturation prevents a direct
application of linear discrepancy bounds. In the sparse setting, however,
each item participates in only a bounded number of feature occurrences.
We show how to exploit that sparsity without any dependence on the size
of the instance.

Here is the exact model. Let $n,m,t$ be positive integers. For each
$i\in[m]=\{1,\ldots,m\}$, let $U_i$ be a finite feature-index set. A feature
occurrence is a tagged pair $e=(i,u)$, with $u\in U_i$, having a nonempty
incidence set $S_e\subseteq[n]$. Write
$\calE=\{(i,u):i\in[m],\ u\in U_i\}$, and suppose
\begin{equation}
 d(j):=\sum_{e\in\calE}\ind[j\in S_e]\le t
 \qquad (j\in[n]).
 \label{eq:sparsity}
\end{equation}
The corresponding coverage functions are
\begin{equation}
 f_i(T)=\sum_{u\in U_i}\ind[S_{(i,u)}\cap T\ne\varnothing]
 \qquad (T\subseteq[n]).
 \label{eq:coverage}
\end{equation}
Equal incidence sets at different tags remain distinct occurrences in
both sums. Empty features can be removed. The indices $U_i$ may be empty.
For a coloring $\chi:[n]\to\Zk=\{0,\ldots,k-1\}$, put
$T_c=\chi^{-1}(c)$ and define
\[
 \disc(\chi)=\max_{i\in[m]}\max_{c,c'\in\Zk}
          |f_i(T_c)-f_i(T_{c'})|.
\]
Empty color classes are allowed.

\begin{theorem}
\label{thm:main}
For every family satisfying \eqref{eq:sparsity} and every integer $k\ge2$,
there is a coloring $\chi:[n]\to\Zk$ such that
\begin{equation}
 \disc(\chi)\le 4kt(kt-1)^2.
 \label{eq:bound}
\end{equation}
For an input that explicitly lists item slots, function and feature tags,
and all incidences, such a coloring can be computed deterministically in
time polynomial in the input length, with $k,t$ encoded in binary.
\end{theorem}

In particular, $P(t,k)=4kt(kt-1)^2$ is one fixed polynomial valid for all
$n,m,t,k$ in this model. No bound on feature sizes is assumed.

\paragraph{Closest prior work and the change in parameters.}
Avila, Rohwedder and Wennmann~\cite{ARW26}, Theorem~3, prove a randomized
polynomial-time bound
$O(\sqrt{t^3k}\log(ntk))$ for sparse coverage functions. Their
Section~1.3 discusses avoiding the logarithmic size factor. Our theorem
gives a positive answer to the qualitative polynomial-existence question
suggested there: every fixed $t,k$ has a bound uniform over all instance
sizes. The price is the coarser $O(t^3k^3)$ dependence. We do not establish
$O(\sqrt{t^3k})$, nor do we claim a pointwise improvement over their bound
for all parameter ranges.

The sparsity convention is essential. We use the covering-universe
description in \cite[Section~1.1]{ARW26}: each item covers at most $t$
feature elements across all functions. Their dual-set notation in
Definition~4 uses a union of collections, which could be read as
deduplicating equal incidence subsets. Our explicit tags avoid that
ambiguity. The theorem concerns occurrence-counted sparsity as in
\eqref{eq:sparsity}; it does not assert a bound under literal
deduplication of repeated subsets.

\paragraph{The technical idea.}
The proof has two requirements that must hold in the same choice space.
Large features must keep all colors throughout rounding, while small
features must admit a controlled linearization. Freezing arbitrary colors
on representatives can bias small features. Rounding witness items
separately can destroy their rainbow property. Instead, we keep each
witness as one categorical variable whose states are cyclic shifts.
Every state protects its large feature, and the uniform distribution on
states treats colors symmetrically. Within a suitable interaction class,
a small feature meets at most one moving block, although it may contain
several items of that block. Its expectation is still affine in the
categorical variable. The crucial column estimate counts the incidences
of those items, giving mass at most $kt$ across \emph{all} functions and
colors.

Independent interaction classes and affine restricted coverage
expectations already appear in \cite[Section~5, Lemma~23]{ARW26}. The
multiple-choice rounding theorem used here is also known, due to Doerr
and Srivastav~\cite{DS03}; its exact rational formulation is stated by
Gla{\ss}er, Reitwie{\ss}ner and Witek~\cite[Theorem~2]{GRW11}. Hall
matching, greedy graph coloring, and that rounding theorem are prior
tools. The structural step here is the choice of disjoint rainbow
witnesses together with categorical coverage extensions that preserve
both requirements and need only a size-independent number of stages.

\paragraph{Broader scope.}
Hollender, Manurangsi, Meka and Suksompong~\cite{HMMS25} and Dupr{\'e}
la Tour and Fujii~\cite{DF25} study discrepancy beyond additive functions.
The general influence-sparse Lipschitz question in
\cite[Definition~13 and Question~14]{HMMS25} has a broader function class
and a different sparsity measure. The present theorem does not resolve
it. Our hypotheses also exclude arbitrary real feature weights and a
bound only on the number of influenced functions.

\section{Multiple-choice rounding}
\label{sec:rounding}

We give the rounding argument in full to specify exactly the prior result
on which the coverage reduction relies. For a finite matrix $A$, write
$\|A\|_1=\max_a\sum_b|A_{b,a}|$ for its maximum absolute column sum.

\begin{lemma}[Multiple-choice Beck--Fiala rounding]
\label{lem:round}
Partition a finite set of coordinates into nonempty groups. Suppose
$p\ge0$ sums to one in every group, and $A$ is a real matrix with these
coordinates as columns and $\|A\|_1\le D$. There is a vector $y$ with
exactly one entry equal to one in each group and all other entries zero
such that
\begin{equation}
 \|A(y-p)\|_\infty\le 2D.
 \label{eq:round}
\end{equation}
For rational data and equally sized groups, a suitable $y$ is computable
in polynomial time.
\end{lemma}

The rational equal-group statement and its algorithmic claim are
\cite[Theorem~2]{GRW11}, attributed there to \cite{DS03}. The following
floating-coordinate proof establishes the existence statement, also for
unequal group sizes; that proof is not claimed as a contribution.

\begin{proof}
Start with $x=p$. Fix permanently every coordinate already equal to zero
or one. Initially retain every row equation $A_bx=A_bp$. At any stage,
let $W$ be the strictly fractional coordinates and put $V=|W|$. If
$V=0$, the process is finished. A group meeting $W$ has at least two
coordinates in $W$, since a nonnegative sum-one group cannot have exactly
one fractional coordinate and all others integral. Thus the number $G$
of such groups satisfies $G\le V/2$.

Drop each retained row $b$ whose residual mass
$\sum_{a\in W}|A_{b,a}|$ is at most $2D$. If $D>0$, the number $R$ of
rows still retained satisfies
\[
 2DR < \sum_{\text{retained }b}\sum_{a\in W}|A_{b,a}|
       \le DV,
\]
and hence $R<V/2$. If $D=0$, all matrix entries vanish and all rows have
been dropped, so again $R+G<V$. There is therefore a nonzero direction
supported on $W$ preserving the retained row equations and the $G$
group-sum equations. Move in that direction until the first coordinate
reaches zero or one, and fix all new boundary coordinates. A positive
finite move is possible because every coordinate of $W$ starts strictly
inside $(0,1)$. Repeat. At least one additional coordinate is fixed at
each move, so there are finitely many moves and the result is an integral
group choice $y$.

A retained row is exact until it is dropped. If row $b$ is dropped at
the point $x^*$ with fractional coordinate set $W^*$, only coordinates
in $W^*$ can move later. Consequently
\[
 |A_b(y-p)|
 =\left|\sum_{a\in W^*}A_{b,a}(y_a-x^*_a)\right|
 \le \sum_{a\in W^*}|A_{b,a}|\le 2D.
\]
A row never dropped remains exact. The error bound charges a row only at
its single deletion, rather than at every move, proving \eqref{eq:round}.
The rational running-time claim can be obtained from the cited theorem;
Section~\ref{sec:algorithm} also describes a rational LP implementation.
\end{proof}

\section{Protecting large features and separating small ones}
\label{sec:blocks}

Set $q=kt\ge2$. Call an occurrence \emph{large} if $|S_e|\ge q$, and
\emph{small} otherwise.

\begin{lemma}
\label{lem:witness}
There are pairwise disjoint sets $B_e\subseteq S_e$, each of size $k$,
one for every large occurrence $e$.
\end{lemma}
\begin{proof}
For any subfamily $Q$ of large tagged occurrences, count incidences using
\eqref{eq:sparsity}:
\[
 kt|Q|\le\sum_{e\in Q}|S_e|
         \le t\left|\bigcup_{e\in Q}S_e\right|.
\]
Thus $|\bigcup_{e\in Q}S_e|\ge k|Q|$. Form a bipartite graph with $k$
demand copies of every large occurrence on the left and the items on the
right; each copy of $e$ is adjacent to $S_e$. For any subset $X$ of
copies, let $Q$ be the occurrences represented in $X$. Then
$|N(X)|=|\bigcup_{e\in Q}S_e|\ge k|Q|\ge|X|$. Hall's theorem gives a
matching saturating all copies. Their matched items are the desired
witness sets. The argument counts tags even when incidence sets coincide.
\end{proof}

Add one singleton block for each unmatched item. The witnesses and these
singletons form a partition $\calB$ of $[n]$, with block sizes $k$ or one.
Assign distinct base colors $a(j)\in\Zk$ to the $k$ items of each witness,
and base color zero to each singleton. Block $B$ has $k$ allowed states:
in state $r\in\Zk$, its items receive colors
\begin{equation}
 \chi(j)=a(j)+r\pmod k \qquad(j\in B).
 \label{eq:states}
\end{equation}
In every state each witness contains every color. Every large occurrence
therefore contributes exactly one to each color's coverage, for every
choice of states of all blocks. It will remain a color-independent
constant throughout the proof.

Define a simple graph on $\calB$, joining distinct blocks whenever a small
occurrence meets both.

\begin{lemma}
\label{lem:classes}
The block graph has a proper coloring with
$H\le(q-1)^2$ classes $\calC_1,\ldots,\calC_H$. Each small occurrence
meets at most one block in any one class.
\end{lemma}
\begin{proof}
A block $B$ meets at most $|B|t\le q$ tagged occurrences. A small
occurrence has at most $q-1$ items, so if it meets $B$ it meets at most
$q-2$ other blocks. Hence the degree of $B$ is at most $q(q-2)$,
where counting neighbors more than once only increases the estimate.
Greedy coloring uses at most $q(q-2)+1=(q-1)^2$ classes. Blocks in one
class cannot both meet a small occurrence, by definition of adjacency.
When $q=2$, every small occurrence is a singleton and the graph is
edgeless, as the formula requires.
\end{proof}

\section{Affine coverage on a class of blocks}
\label{sec:affine}

Give every block $B$ a probability vector
$p_B=(p_{B,r})_{r\in\Zk}$ on its allowed states. Choose states
independently \emph{between blocks}, retaining all dependencies between
items of the same block. For the resulting product law define
\[
 \Phiic(p)=\mathbb{E}[f_i(T_c)].
\]
Initially let every $p_{B,r}=1/k$. Adding the same shift to every block
rotates all colors and preserves this product law. Thus there is a common
value $\mu_i$ with
\begin{equation}
 \Phiic(p^{(0)})=\mu_i \qquad(c\in\Zk).
 \label{eq:symmetry}
\end{equation}
This symmetry does not require independent colors within a block.

Fix a class $\calC$ and hold all probabilities outside $\calC$ fixed.
For a small occurrence $e$ meeting $B\in\calC$, define
\begin{align*}
 z_{e,c}(r)&=\ind[\text{some }j\in S_e\cap B
                         \text{ has }a(j)+r=c\pmod k],\\
 v_{e,c}&=\Pr[\text{no block outside }\calC
                         \text{ gives color }c\text{ to }S_e].
\end{align*}
The occurrence meets at most one block of $\calC$, by
Lemma~\ref{lem:classes}. Independence of that block and the outside
blocks therefore gives the exact hit probability
\begin{equation}
 \Pr[S_e\cap T_c\ne\varnothing]
    =1-v_{e,c}+v_{e,c}\sum_{r\in\Zk}p_{B,r}z_{e,c}(r).
 \label{eq:feature-affine}
\end{equation}
If $e$ meets no class block, its contribution is constant. Summing
\eqref{eq:feature-affine} over small occurrences and adding the constant
large-feature contributions gives an affine map
\begin{equation}
 (\Phiic(p))_{i,c}=b+A p_{\calC},\qquad
 A_{(i,c),(B,r)}=
 \sum_{\substack{e=(i,u)\text{ small}\\S_e\cap B\ne\varnothing}}
                    v_{e,c}z_{e,c}(r).
 \label{eq:matrix}
\end{equation}
The vector $b$ and matrix $A$ depend on the fixed outside laws.

\begin{lemma}
\label{lem:column}
The matrix in \eqref{eq:matrix} has nonnegative entries and
$\|A\|_1\le q$.
\end{lemma}
\begin{proof}
For fixed $B,r$, all items of $B$ receive distinct colors. Consequently
$\sum_{c\in\Zk}z_{e,c}(r)=|S_e\cap B|$, including when $e$ contains
several items of $B$. Since $0\le v_{e,c}\le1$,
\begin{align}
 \sum_{i,c} A_{(i,c),(B,r)}
 &\le\sum_{\substack{e\text{ small}\\S_e\cap B\ne\varnothing}}
                       |S_e\cap B| \notag\\
 &=\sum_{j\in B}d_{\rm small}(j)
 \le |B|t\le q,
 \label{eq:column}
\end{align}
where $d_{\rm small}(j)$ counts small tagged occurrences containing $j$.
All terms in \eqref{eq:matrix} are nonnegative. This proves the claimed
absolute column bound. In particular, multiplicities are included and no
extra factor equal to the number of colors or functions is introduced.
\end{proof}

\section{Rounding the classes}
\label{sec:finish}

\begin{proof}[Proof of the existence assertion in Theorem~\ref{thm:main}]
Starting from $p^{(0)}$, process the classes in order. In class
$\calC_h$, hold all outside laws fixed and apply Lemma~\ref{lem:round}
to \eqref{eq:matrix}, with one $k$-coordinate group per block and $D=q$.
Make the states in this class deterministic using the integral group
choices returned by the lemma. If $p^{(h)}$ is the full collection of
block laws after this step, then for every $i,c$,
\begin{equation}
 |\Phiic(p^{(h)})-\Phiic(p^{(h-1)})|\le2q.
 \label{eq:stage}
\end{equation}
At each step the matrix is recomputed using the actual outside laws.
Earlier classes are point masses, and later classes are still uniform.
Both are legitimate marginals of the same product-law definition of
$\Phi$. Thus \eqref{eq:stage} compares consecutive evaluations of one
functional; it does not require the matrices of different classes to
coincide.

After $H$ stages every block law is a point mass. The final $\Phiic$ is
the actual coverage $f_i(T_c)$ of the selected coloring. Telescoping
\eqref{eq:stage} and using \eqref{eq:symmetry} gives
$|f_i(T_c)-\mu_i|\le2qH$. Subtracting the same $\mu_i$ for two colors
and using Lemma~\ref{lem:classes} yields
\[
 |f_i(T_c)-f_i(T_{c'})|\le4qH\le4kt(kt-1)^2.
\]
This holds simultaneously for all functions and pairs of colors.
\end{proof}

\section{Deterministic construction and encoding length}
\label{sec:algorithm}

We justify the time claim in the explicit input model of
Theorem~\ref{thm:main}. In particular, a large binary value of $k$ must not
force an exponentially expanded state space.

If $k\ge n$, assign distinct colors to the items. A nonempty color class
is a singleton, whose value for any $f_i$ is at most $d(j)\le t$; empty
classes have value zero. The discrepancy is therefore at most $t$, which
implies \eqref{eq:bound}. The output labels have polynomial length.
Hence assume $k<n$. Replace the supplied $t$, for implementation purposes,
by $d=\max(1,\max_j d(j))\le t$. The function
$x\mapsto4kx(kx-1)^2$ is increasing for $x\ge1$ and $k\ge2$, so the
resulting guarantee implies \eqref{eq:bound}. Both $k$ and $d$ are now
polynomially bounded by the explicit input length. Features and functions
with no incidence can be discarded.

Use the threshold $kd$ in Lemma~\ref{lem:witness}. The matching graph has
at most $n$ demand copies: the proof's counting inequality, applied to
all large occurrences, gives this directly. Its edge list is polynomial
in the incidence input because $k<n$. Finite bipartite matching, block
graph construction, and greedy coloring are polynomial time. The actual
number of nonempty classes is at most the number of blocks, hence at
most $n$, regardless of the numerical upper bound $(kd-1)^2$.

It remains to control the rational matrices used for rounding. At the
start of a class, every outside block is either uniform or deterministic.
For a uniform outside block $B$, the probability of missing a fixed
color on $S_e\cap B$ equals
\[
 \frac{k-|S_e\cap B|}{k},
\]
because its base colors are distinct. A deterministic outside block has
miss probability zero or one. Therefore $v_{e,c}$ is a product of at most
$n$ rational factors with denominator $k$. If $M=|\calE|$, every
coefficient of \eqref{eq:matrix} has a representation with denominator
dividing $k^n$ and bit length $O(n\log k+\log(M+1))$. The initial target
for the current uniform class has denominator dividing $k^{n+1}$.
There are at most $mk$ rows and $nk$ columns. The rational matrix and
target thus have polynomial encoding length and can be computed exactly.
The polynomial-time equal-group version of Lemma~\ref{lem:round} applies.

For completeness, one may implement the displayed rounding proof without
assuming a bit bound for arbitrary repeated nullspace moves. After
dropping light rows, consider the bounded rational polytope defined by
the remaining row equations, group sums, bounds $0\le x_a\le1$, and all
previously fixed coordinates. It is nonempty. Since $R+G<V$, a point
strictly inside the bounds of all $V$ remaining coordinates cannot be a
vertex: a nonzero homogeneous direction would permit movement in both
directions. Thus any vertex fixes at least one further coordinate.
Obtain a vertex by sequentially optimizing coordinates and restricting
to the corresponding optimal faces. Polynomial-time rational linear
programming suffices. Each retained row has its original target $A_bp$
for this class, adjusted only by original coefficients of coordinates
fixed at zero or one. The constraint coefficients and right sides
therefore keep polynomial bit length, and determinant bounds give the
same property for vertex coordinates. There are at most $nk$ coordinate
fixings per class. This yields a deterministic polynomial-time
implementation and completes Theorem~\ref{thm:main}.

\section{Limits and verification}

The result gives existence and an explicit-input algorithm for all
parameters in \eqref{eq:sparsity}--\eqref{eq:coverage}. It does not optimize
the $t,k$ exponents. A single item covering $t$ distinct features has
discrepancy $t$ for every $k\ge2$, the lower-bound example also discussed
in \cite[Section~1.3]{ARW26}; a substantial parameter gap remains. The
categorical construction does not by itself address arbitrary weighted
coverage or arbitrary influence-sparse Lipschitz functions.

The proof uses no computational search. Supplementary exact-rational
diagnostics tested the categorical identities, witness conditions, and
grouped rounding on finite instances. A separate model context also
reconstructed the argument and checked finite instances using a separately
written exact-arithmetic program. These checks do not establish the
infinite theorem, replace the proof, certify priority, or constitute
external peer review. The comparison with the cited primary results is
limited to the literature inspected.

\paragraph{AI methodology.}
OpenAI Codex provided substantial assistance with the proof development,
counterexample checks, exact-arithmetic code, literature comparison, and
the writing of this manuscript; this assistance materially affected all
sections. A separate model context assisted with proof and source checks.
No claim of human verification or formal proof certification is made.

\clearpage
\bibliographystyle{alpha}
\bibliography{references}
\end{document}
